% chapters/14/praktikum.tex
\begin{praktikumbox}
Praktikum ini mengikuti Tugas Praktikum 11 pada RPS: mengimpor aset, membuat
animasi diam dan berlari, menampilkan skor pada antarmuka, dan membandingkan
scene sebelum dan sesudah penggantian aset.

\begin{langkah}
  \item \textbf{Impor asetnya.} Unduh sekurang-kurangnya satu aset pemain yang
        memuat animasi diam dan berlari, beserta aset untuk latar dan
        rintangan. Catat sumber dan lisensi setiap aset sejak awal.

  \item \textbf{Atur pengaturan impornya.} Atur \istilah{Pixels Per Unit} pada
        seluruh aset agar nilainya sama, lalu tentukan tinggi pemain yang
        dikehendaki dan hitunglah nilai yang diperlukan.

  \item \textbf{Potong \istilah{sprite sheet}-nya.} Apabila asetnya berupa
        \istilah{sprite sheet}, potonglah menjadi bingkai-bingkai menurut ukuran
        yang seragam.

  \item \textbf{Susun \istilah{Animator}.} Buat keadaan diam dan berlari,
        tambahkan parameter kecepatan dan keadaan menyentuh tanah, lalu
        gambarlah transisinya.

  \item \textbf{Tempelkan skrip animasinya.} Ketik ulang
        \code{ch14_kontrol_animasi.cs} dari Daftar~\ref{lst:ch14-animasi}, lalu
        hubungkan rujukan \istilah{Animator} dan \istilah{Rigidbody} pada
        \istilah{Inspector}. Jalankan scene-nya dan pastikan animasinya
        berganti menurut gerak pemain.

  \item \textbf{Susun antarmuka pengguna.} Buat \istilah{Canvas} berisi teks
        skor dan teks nyawa, lalu tempelkan skrip
        \code{ch14_tampilan_hud.cs} dari Daftar~\ref{lst:14}. Hubungkan
        rujukan teksnya, lalu uji dengan menambahkan perintah pemanggil
        \code{TambahSkor} ketika pemain mengambil benda.

  \item \textbf{Bandingkan sebelum dan sesudah.} Simpan tangkapan layar scene
        sebelum dan sesudah penggantian aset, lalu periksa apakah mekaniknya
        berubah. Perbaiki bentuk batas benda apabila ukurannya berubah.

  \item \textbf{Tulis catatan praktikum} yang memuat daftar aset beserta
        sumber dan lisensinya, pengaturan impor yang dipakai, tangkapan layar
        sebelum dan sesudah, serta masalah yang ditemui dan cara
        mengatasinya.

  \item \textbf{Simpan pekerjaan} ke repositori Git dengan pesan penyimpanan
        yang menerangkan isinya.
\end{langkah}
\end{praktikumbox}

\begin{table}[htbp]
\centering
\caption{Lembar kerja Praktikum 11}
\label{tab:lembar-praktikum-11}
{\setlength{\tabcolsep}{4pt}\footnotesize%
\begin{tabularx}{\textwidth}{@{}L{4.2cm} C{2.2cm} Y@{}}
\toprule
\headrow \textbf{Nama aset} & \textbf{Lisensi} & \textbf{Sumber dan keperluan atribusi} \\
\midrule
& & \\[1.6em]
\midrule
& & \\[1.6em]
\midrule
& & \\[1.6em]
\midrule
\addlinespace
\multicolumn{3}{@{}p{\textwidth}@{}}{\footnotesize\itshape
Nilai \istilah{Pixels Per Unit} yang dipakai: \dotfill} \\[1.4em]
\addlinespace
\multicolumn{3}{@{}p{\textwidth}@{}}{\footnotesize\itshape
Perubahan mekanik yang timbul sesudah penggantian aset: \dotfill} \\[1.4em]
\addlinespace
\multicolumn{3}{@{}p{\textwidth}@{}}{\footnotesize\itshape
Masalah yang ditemui dan cara mengatasinya: \dotfill} \\[1.8em]
\bottomrule
\end{tabularx}}
\end{table}

\begin{catatabox}
Langkah ketujuh adalah langkah yang paling penting bagi indikator \ind{a}.
Mengganti aset memang mudah, tetapi memastikan bahwa mekaniknya tidak berubah
menuntut pemeriksaan yang sengaja. Bandingkanlah kedua tangkapan layar itu
dengan saksama.
\end{catatabox}
